% GENERATED, edit the French unit or the translation table
\mxnewpage\section*{EV1-A -- Common test}\label{EV1-A}
\addcontentsline{toc}{section}{EV1-A -- Common test}
\textbf{Paper reserved for the teacher}. Chapters 1 to 4.\par Duration: 2 h. No calculator. Marked out of 20. 14 November 2026.
\subsection*{Part 1 -- Fluency (6 points)}
Twelve questions, 0.5 point each. Only one correct answer; no justification required. An incorrect or missing answer scores 0, with no loss of points.
\begin{enumerate}
\item Calculate $5+2\times3$.\par A. $14$\quad B. $11$\quad C. $10$\quad D. $21$
\item Calculate $\frac12+\frac13$.\par A. $\frac{5}{6}$\quad B. $\frac{2}{5}$\quad C. $\frac{5}{3}$\quad D. $\frac{1}{6}$
\item Calculate $\dfrac{2^5}{2^{3}}$.\par A. $2$\quad B. $1$\quad C. $10$\quad D. $4$
\item A value of 50 increases by 20 \%. Final value?\par A. $10$\quad B. $60$\quad C. $70$\quad D. $40$
\item A value of 80 is obtained after a 20 \% decrease. Initial value?\par A. $100$\quad B. $60$\quad C. $64$\quad D. $96$
\item A 25 \% increase then a 20 \% decrease: overall coefficient?\par A. $\frac{9}{10}$\quad B. $1$\quad C. $\frac{21}{20}$\quad D. $\frac{4}{5}$
\item The solution of $2x+5=0$ is\par A. $10$\quad B. $\frac{-5}{2}$\quad C. $-10$\quad D. $\frac{5}{2}$
\item The slope between $(0;5)$ and $(2;11)$ is\par A. $6$\quad B. $3$\quad C. $2$\quad D. $-3$
\item The mean of $5, 7, 9$ is\par A. $8$\quad B. $21$\quad C. $5$\quad D. $7$
\item A sector of $90^\circ$ represents what fraction of the total?\par A. $\frac{1}{3}$\quad B. $\frac{1}{2}$\quad C. $\frac{1}{4}$\quad D. $\frac{3}{4}$
\item A probability is $\dfrac{2}{5}$. Its complement is\par A. $\frac{3}{5}$\quad B. $\frac{2}{5}$\quad C. $\frac{5}{2}$\quad D. $\frac{1}{5}$
\item For $f(x)=x^2+5$, $f(-2)$ is\par A. $9$\quad B. $7$\quad C. $1$\quad D. $20$
\end{enumerate}
\mxnewpage\subsection*{Part 2 -- Independent exercises (14 points)}
Every step must be justified. Exact results are expected unless stated otherwise.
\mxneed{12}
\subsection*{Exercise 1 -- Decreasing reserve (5 points)}
A tank initially contains 192 L of water. At each step, half of the water still present is removed, with no water added. Let $u_n$ be the volume, in litres, remaining after $n$ steps; index 0 corresponds to the initial state.
\begin{enumerate}
\item Define the sequence with its first term and common ratio. \hfill(1 pt)
\item Calculate the volumes for indices 0 to 4 and their sum. \hfill(2 pts)
\item Determine the first index where the volume is strictly less than 10 L; describe a threshold loop. \hfill(2 pts)
\end{enumerate}
\mxneed{12}
\subsection*{Exercise 2 -- Parameter and tangent (5 points)}
For a fixed real number $m$, define on $\R$ the function $f_m(x)=x^2-2mx+4$. The number $m$ is a parameter. The first two questions concern all its values; the last concerns only $m=1$.
\begin{enumerate}
\item Calculate the discriminant of the polynomial $f_m$ in terms of $m$. \hfill(1 pt)
\item Determine, according to $m$, the number of real solutions of $f_m(x)=0$. \hfill(2 pts)
\item Now fix $m=1$. Calculate the rate of change between 2 and $2+h$, for $h\ne0$, then the equation of the tangent at the point with x-coordinate 2. \hfill(2 pts)
\end{enumerate}
\mxneed{12}
\subsection*{Exercise 3 -- Locus via the midpoint (4 points)}
Two fixed points $A$ and $B$ of the plane are 8 units apart. Let $I$ be the midpoint of the segment $[AB]$ and $M$ a variable point of the plane. We seek to describe sets of points starting from the scalar product $\vect{MA}\cdot\vect{MB}$.
\begin{enumerate}
\item Express $\vect{MA}\cdot\vect{MB}$ in terms of $MI$, justifying the identity used. \hfill(1 pt)
\item Describe the set of points $M$ for which this product equals 9. \hfill(1 pt)
\item Describe in turn the sets obtained when the product equals $-16$, then $-20$. \hfill(2 pts)
\end{enumerate}
